% This file is generated from the corresponding Obsidian note.
% Source: Teoricos/GS - Teo22.md
\chapter{Complejo de de Rham y pullbacks}
\label{ch:teo-22}
\section{Complejo de de Rham}
\begin{bookdefinition}{Complejo de Rham}
Es el complejo de cocadenas

\[
0\xrightarrow{d}\Omega^0(M)\xrightarrow{d}\Omega^1(M)\xrightarrow{d}\cdots\xrightarrow{d}\Omega^n(M)\xrightarrow{d}0.
\]
Algunas veces se denota por $d_k=d|_{\Omega^k(M)}$, y asi, la condicion $d\circ d=0$, se lee como $\operatorname{Im}d_k\subset\ker d_{k+1}$. Se define el $k$-esimo grupo de cohomologia de de Rham, denotado por

\[
H^k_{dR}(M):=\ker d_k/\operatorname{Im}d_{k-1},
\]
donde

\[
Z^k(M)=\ker d_k,
\]
\[
B^k(M)=\operatorname{Im}d_{k-1}.
\]
\end{bookdefinition}
\begin{bookexample}{}
Por ejemplo:

\[
H^0_{dR}(M)=\ker d_0/\operatorname{Im}d_{-1}=\ker d_0=\{f\in C^\infty(M):df=0\}.
\]
$f$ es constante en cada componente conexa de $M$, entonces $H^0_{dR}(M)\cong\mathbb R^{\#\operatorname{comp}}$. Cuando $M$ es compacta, se puede probar que $\dim H^k_{dR}(M)$ es finito y se llama el $k$-esimo numero de Betti.

\end{bookexample}
\begin{bookdefinition}{Formas cerradas y exactas}
Una forma diferencial $\theta$ se dice cerrada si $d\theta=0$, y se dice exacta si $\theta=d\omega$, para alguna forma diferenciable $\omega$.

\end{bookdefinition}
\begin{bookexample}{}
Sea $B^n=\{x\in\mathbb R^n:\|x\|<1\}$. Se puede probar que

\[
H^k(B^n)=0,\qquad k\ge1,
\]
y ya sabemos $H^0(B^n)=\mathbb R$. (Se llama lema de Poincare)

\end{bookexample}
\begin{bookexample}{}
Se puede probar que si $M$ es simplemente conexa, entonces $H^1_{dR}(M)=0$. En general

\[
H^1_{dR}(M)\cong \operatorname{Hom}(\pi_1(M),\mathbb R).
\]
\[
\pi_1(S^1)=\mathbb Z,\qquad \operatorname{Hom}(\pi_1(M),\mathbb R)\cong\mathbb R.
\]
\end{bookexample}
\section{Pullback de formas}
\begin{bookdefinition}{Pullback de formas}
Sea $F:M\to N$ suave y $\omega\in\Omega^k(N)$. Se define el pullback $F^*\omega$ de $\omega$ como la $k$-forma (se prueba en el siguiente enunciado que es suave)

\[
(F^*\omega)_p(v_1,\ldots,v_k)=\omega_{F(p)}((dF)_pv_1,\ldots,(dF)_pv_k),
\]
para todo $v_1,\ldots,v_k\in T_pM$.

\end{bookdefinition}
\begin{bookproposition}{Propiedades del pullback}
Sea $F:M\to N$ suave.

\begin{itemize}
\item (a) Si $\omega\in\Omega^k(N)$, entonces $F^*\omega\in\Omega^k(M)$, y $F^*:\Omega^k(N)\to\Omega^k(M)$ es transf. lineal.
\item (b) $F^*(\omega\wedge\theta)=(F^*\omega)\wedge(F^*\theta)$.
\item (c) $F^*(d_N\omega)=d_MF^*\omega$.
\end{itemize}
\end{bookproposition}
\begin{bookproof}{Ejercicio}
\begin{itemize}
\item $(a)$
\end{itemize}
\begin{enumerate}
\item Probemos que la sección $p\mapsto(F^*\omega)_p$ es suave. Fijemos $p\in M$ y tomemos cartas $(U,x)$ de $M$ alrededor de $p$ y $(V,y)$ de $N$ alrededor de $F(p)$, achicando $U$ de modo que $F(U)\subseteq V$. En $V$ escribimos
\end{enumerate}
\[
\omega=\sum_I a_I\,dy^{i_1}\wedge\cdots\wedge dy^{i_k},
\]
donde cada $a_I$ es suave. Afirmamos que sobre $U$ se cumple

\[
F^*\omega=\sum_I(a_I\circ F)\,d(y^{i_1}\circ F)\wedge\cdots\wedge d(y^{i_k}\circ F).
\]
\begin{enumerate}[start=2]
\item Para verificar esta igualdad, sean $q\in U$ y $v_1,\ldots,v_k\in T_qM$. Por definición del pullback,
\end{enumerate}
\[
\begin{aligned}(F^*\omega)_q(v_1,\ldots,v_k)&=\omega_{F(q)}\bigl((dF)_q(v_1),\ldots,(dF)_q(v_k)\bigr)\\&=\sum_I a_I(F(q))(dy^{i_1}\wedge\cdots\wedge dy^{i_k})_{F(q)}\bigl((dF)_q(v_1),\ldots,(dF)_q(v_k)\bigr).\end{aligned}
\]
\begin{enumerate}[start=3]
\item Por definición del producto wedge de $1$-formas,
\end{enumerate}
\[
(dy^{i_1}\wedge\cdots\wedge dy^{i_k})_{F(q)}\bigl((dF)_q(v_1),\ldots,(dF)_q(v_k)\bigr)=\det\left(dy^{i_r}_{F(q)}\bigl((dF)_q(v_s)\bigr)\right)_{r,s}.
\]
\begin{enumerate}[start=4]
\item Por la regla de la cadena, $dy^{i_r}_{F(q)}\circ(dF)_q=d(y^{i_r}\circ F)_q$. Por tanto,
\end{enumerate}
\[
\begin{aligned}\det\left(dy^{i_r}_{F(q)}\bigl((dF)_q(v_s)\bigr)\right)_{r,s}&=\det\left(d(y^{i_r}\circ F)_q(v_s)\right)_{r,s}\\&=\bigl(d(y^{i_1}\circ F)\wedge\cdots\wedge d(y^{i_k}\circ F)\bigr)_q(v_1,\ldots,v_k).\end{aligned}
\]
\begin{enumerate}[start=5]
\item Sustituyendo, obtenemos
\end{enumerate}
\[
\begin{aligned}(F^*\omega)_q(v_1,\ldots,v_k)&=\sum_I(a_I\circ F)(q)\bigl(d(y^{i_1}\circ F)\wedge\cdots\wedge d(y^{i_k}\circ F)\bigr)_q(v_1,\ldots,v_k)\\&=\left(\sum_I(a_I\circ F)\,d(y^{i_1}\circ F)\wedge\cdots\wedge d(y^{i_k}\circ F)\right)_q(v_1,\ldots,v_k).\end{aligned}
\]
\begin{enumerate}[start=6]
\item Como esto vale para todo $q\in U$ y todos $v_1,\ldots,v_k\in T_qM$, concluimos que
\end{enumerate}
\[
F^*\omega=\sum_I(a_I\circ F)\,d(y^{i_1}\circ F)\wedge\cdots\wedge d(y^{i_k}\circ F)
\]
sobre $U$.

\begin{enumerate}[start=7]
\item Cada función $a_I\circ F$ es suave y cada $d(y^{i_r}\circ F)$ es una $1$-forma suave. Por lo tanto, el lado derecho es una $k$-forma suave sobre $U$. Como $p$ era arbitrario, $F^*\omega\in\Omega^k(M)$.
\item Finalmente, sean $\omega,\eta\in\Omega^k(N)$ y $a,b\in\mathbb R$. Para todo $p\in M$ y $v_1,\ldots,v_k\in T_pM$,
\end{enumerate}
\[
\begin{aligned}\bigl(F^*(a\omega+b\eta)\bigr)_p(v_1,\ldots,v_k)&=(a\omega+b\eta)_{F(p)}\bigl((dF)_p(v_1),\ldots,(dF)_p(v_k)\bigr)\\&=a\,\omega_{F(p)}\bigl((dF)_p(v_1),\ldots,(dF)_p(v_k)\bigr)\\&\quad+b\,\eta_{F(p)}\bigl((dF)_p(v_1),\ldots,(dF)_p(v_k)\bigr)\\&=\bigl(aF^*\omega+bF^*\eta\bigr)_p(v_1,\ldots,v_k).\end{aligned}
\]
Luego $F^*(a\omega+b\eta)=aF^*\omega+bF^*\eta$, por lo que $F^*$ es lineal.

\begin{itemize}
\item $(b)$
\end{itemize}
\begin{enumerate}
\item Para probar que $F^*(\omega\wedge\theta)=(F^*\omega)\wedge(F^*\theta)$, evaluamos ambos lados en $p\in M$ y en vectores $v_1,\ldots,v_{k+\ell}\in T_pM$.
\item Por definición de pullback,
\end{enumerate}
\[
(F^*(\omega\wedge\theta))_p(v_1,\ldots,v_{k+\ell})=(\omega\wedge\theta)_{F(p)}((dF)_pv_1,\ldots,(dF)_pv_{k+\ell}).
\]
\begin{enumerate}[start=3]
\item Usando la definición de producto cuña,
\end{enumerate}
\[
(\omega\wedge\theta)_{F(p)}((dF)_pv_1,\ldots,(dF)_pv_{k+\ell})=\frac{1}{k!\ell!}\sum_{\pi\in S_{k+\ell}}\operatorname{sgn}(\pi)\,\omega_{F(p)}((dF)_pv_{\pi(1)},\ldots,(dF)_pv_{\pi(k)})\,\theta_{F(p)}((dF)_pv_{\pi(k+1)},\ldots,(dF)_pv_{\pi(k+\ell)}).
\]
\begin{enumerate}[start=4]
\item Pero, otra vez por definición de pullback,
\end{enumerate}
\[
\omega_{F(p)}((dF)_pv_{\pi(1)},\ldots,(dF)_pv_{\pi(k)})=(F^*\omega)_p(v_{\pi(1)},\ldots,v_{\pi(k)}),
\]
y

\[
\theta_{F(p)}((dF)_pv_{\pi(k+1)},\ldots,(dF)_pv_{\pi(k+\ell)})=(F^*\theta)_p(v_{\pi(k+1)},\ldots,v_{\pi(k+\ell)}).
\]
\begin{enumerate}[start=5]
\item Entonces
\end{enumerate}
\[
(F^*(\omega\wedge\theta))_p(v_1,\ldots,v_{k+\ell})=((F^*\omega)\wedge(F^*\theta))_p(v_1,\ldots,v_{k+\ell}).
\]
Como esto vale para todo $p$ y todo $v_1,\ldots,v_{k+\ell}$, queda

\[
F^*(\omega\wedge\theta)=(F^*\omega)\wedge(F^*\theta).
\]
\begin{itemize}
\item $(c)$
\end{itemize}
\begin{enumerate}
\item Por linealidad de $F^*$ y de la derivada exterior, es suficiente con verificar la igualdad en formas expresadas sobre un abierto coordenado $U$.
\item Recordemos que por definición de pullback de $0$-formas,
\item 
\end{enumerate}
\[
F^*(f)=f\circ F.
\]
\begin{enumerate}[start=4]
\item Así pues, dado $p\in M$ y $v\in T_pM$,
\end{enumerate}
\[
(F^*(df))_pv=(df)_{F(p)}((dF)_pv)=d(f\circ F)_pv.
\]
\begin{enumerate}[start=5]
\item Si $\theta$ es una $k$-forma, como la definición de la derivada exterior es local y no depende de la carta, fijemos un $p\in M$ y una carta $(U,\varphi=(x_1,\ldots,x_n))$ de $F(p)$:
\end{enumerate}
\[
\theta|_U=\sum_{i_1<\cdots<i_k}a_{i_1\ldots i_k}\,dx_{i_1}\wedge\cdots\wedge dx_{i_k}=\sum_Ia_I\,dx_I.
\]
\begin{enumerate}[start=6]
\item Entonces
\end{enumerate}
\[
(d_N\theta)|_U=\sum_I(da_I)\wedge dx_I.
\]
\begin{enumerate}[start=7]
\item Por tanto
\end{enumerate}
\[
\begin{aligned}(F^*(d_N\theta))_p&=\left(F^*\left(\sum_I(da_I)\wedge dx_I\right)\right)_p\\&=\sum_I(F^*(da_I))_p\wedge(F^*(dx_I))_p\\&=\sum_{i_1<\cdots<i_k}d(a_I\circ F)_p\wedge d(x_{i_1}\circ F)_p\wedge\cdots\wedge d(x_{i_k}\circ F)_p.\end{aligned}
\]
En la última igualdad usamos el paso 3.

\begin{enumerate}[start=8]
\item Por otro lado
\end{enumerate}
\[
\begin{aligned}(d_M(F^*\theta))_p&=\left(d_M\left(F^*\left(\sum_Ia_I\,dx_I\right)\right)\right)_p\\&=\left(d_M\left(\sum_I(F^*a_I)\wedge F^*(dx_I)\right)\right)_p\\&=\left(d_M\left(\sum_{i_1<\cdots<i_k}(a_I\circ F)\,d(x_{i_1}\circ F)\wedge\cdots\wedge d(x_{i_k}\circ F)\right)\right)_p\\&=\sum_{i_1<\cdots<i_k}d(a_I\circ F)_p\wedge d(x_{i_1}\circ F)_p\wedge\cdots\wedge d(x_{i_k}\circ F)_p.\end{aligned}
\]
\begin{enumerate}[start=9]
\item En la última igualdad usamos que $d_Md_M=0$. Comparando con la expresión anterior, concluimos que
\end{enumerate}
\[
(F^*(d_N\theta))_p=(d_M(F^*\theta))_p.
\]
\bookqed
\end{bookproof}
